%% kapa: exam=1 round=34 year=2023 period=1 form=A subject=civil_law subject_ko=민법 seq=11 number=11
%% source: Q-Net 감정평가사 기출문제 (cst00309) artlSeq=5211789 file=raw/1008/5211789/2023년도 제34회 감정평가사 1차시험 1교시.pdf
%% conversion: kordoc 4.24.51 -> tools/gen1.py (auto); review=proofread-sonnet; flags=none
\begin{question}{11}
통정허위표시에 관한 설명으로 옳은 것은? (다툼이 있으면 판례에 따름)
\begin{choices}
  \item 통정허위표시에 의한 급부는 특별한 사정이 없는 한 불법원인급여이다.
  \item 대리인이 대리권의 범위 안에서 현명하여 상대방과 통정허위표시를 한 경우, 본인이 선의라면 특별한 사정이 없는 한 그는 허위표시의 유효를 주장할 수 있다.
  \item 가장행위인 매매계약이 무효라면 은닉행위인 증여계약도 당연히 무효이다.
  \item 통정허위표시의 무효로부터 보호되는 선의의 제3자는 통정허위표시를 알지 못한 것에 대해 과실이 없어야 한다.
  \item 가장매매계약의 매수인과 직접 이해관계를 맺은 제3자가 악의라 하더라도 그와 다시 법률상 이해관계를 맺은 전득자가 선의라면 가장매매계약의 무효로써 전득자에게 대항할 수 없다.
\end{choices}
\end{question}
\answer{5}
